% chapter5_results_analysis.tex
% Chapter 5: Results and Analysis
% AI-Based Healthcare Claims Fraud Detection with Medical Validation and Explainable AI
% B.Tech Computer Engineering Major Project — SPIT, Mumbai

\chapter{Results and Analysis}

\section{Experimental Setup}

All experiments were conducted on a standard personal computing environment running
Microsoft Windows with Python 3.13 (64-bit). The complete data processing and model
training pipeline was implemented in Python and executed using the
\verb|scripts/build_data_store.py| script and the
\verb|scripts/generate_eval_data.py| evaluation script. The Streamlit interactive
dashboard was tested on \verb|localhost:8501|. No external GPU or cloud computing
infrastructure was required for training, as the dataset sizes and chosen model
complexities remained within the computational capacity of commodity hardware with
at least 8\,GB of RAM.

The experimental evaluation is structured into two independent components. The first
component evaluates the NFHS-5 medical consistency screening module, which applies
six rule-based clinical inconsistency rules to the nationwide benchmark dataset of
724,115 records and to the 35,000-record screened sample. The second component
evaluates the Medicare provider-level potential fraud detection module using
supervised machine learning. Both components are evaluated independently; no
record-level linkage between NFHS-5 and Medicare data exists within the system.

\begin{table}[h!]
\centering
\renewcommand{\arraystretch}{1.3}
\begin{tabular}{|l|l|l|}
\hline
\textbf{Component} & \textbf{Dataset} & \textbf{Scale} \\
\hline
NFHS-5 Medical Validation & National Family Health Survey-5 & 724,115 records (nationwide) \\
\hline
NFHS-5 Sample (Dashboard) & NFHS-5 Screened Sample & 35,000 records \\
\hline
Provider-Level ML & CMS SynPUF Medicare Claims & 5,410 providers; 558,211 claims \\
\hline
\end{tabular}
\caption{Datasets Used in Experimental Evaluation}
\label{tab:exp-datasets}
\end{table}

\section{Dataset Description and Statistics}

\subsection{NFHS-5 Dataset}

The National Family Health Survey-5 (NFHS-5) is a large-scale, nationally
representative household survey conducted across India during 2019--2021. It was
administered by the International Institute for Population Sciences (IIPS), Mumbai,
under the Ministry of Health and Family Welfare, Government of India. The dataset
used in this project contains approximately 724,115 female respondent records. Each
record contains self-reported and clinically measured health information, including
diabetes diagnosis status, hypertension diagnosis status, treatment and medication
status, and random blood glucose measurements.

For the purpose of this project, clinically relevant variables related to diabetes,
hypertension, treatment, medication, and glucose level were selected. Records
containing invalid survey codes, missing values for all relevant clinical variables,
or otherwise excluded respondents were filtered prior to the rule evaluation. The
preprocessing step handled approximately 724,115 valid records.

\subsection{CMS SynPUF Medicare Dataset}

The Medicare dataset used in this project is derived from the Centers for Medicare
and Medicaid Services (CMS) Synthetic Public Use File (SynPUF) dataset, which is a
publicly available synthetic representation of Medicare claims data. The dataset
consists of four primary files:

\begin{itemize}

\item \textbf{Train Labels:} Provider-level binary target variable
  \verb|PotentialFraud| (``Yes'' or ``No'') for 5,410 unique providers.

\item \textbf{Inpatient Claims:} Claim records for acute hospital inpatient
  admissions, including admission and discharge dates, diagnosis codes (up to 10
  ICD codes), procedure codes, diagnosis group codes, attending and operating
  physician identifiers, and reimbursement amounts. This file contains a
  substantially larger volume of per-claim records.

\item \textbf{Outpatient Claims:} Claim records for outpatient visits and procedures,
  including start and end dates, diagnosis codes (up to 10 ICD codes), procedure
  codes (up to 6), and reimbursement amounts.

\item \textbf{Beneficiary Data:} Demographic records for Medicare beneficiaries,
  including date of birth, date of death (if applicable), gender, state (FIPS code),
  race, and binary flags for 11 chronic conditions (Alzheimer's Disease, Heart
  Failure, Kidney Disease, Cancer, COPD/Pulmonary Disease, Depression, Diabetes,
  Ischaemic Heart Disease, Osteoporosis, Rheumatoid Arthritis, and Stroke).

\end{itemize}

After combining inpatient and outpatient claims, the integrated claims dataset
contains 558,211 claim records. The provider-level dataset, derived by aggregating
claim-level information into provider-level statistical features, contains 5,410
providers.

\begin{table}[h!]
\centering
\renewcommand{\arraystretch}{1.3}
\begin{tabular}{|l|r|r|}
\hline
\textbf{Class} & \textbf{Providers} & \textbf{Percentage} \\
\hline
No Potential Fraud & 4,904 & 90.65\% \\
\hline
Potential Fraud & 506 & 9.35\% \\
\hline
\textbf{Total} & \textbf{5,410} & \textbf{100.00\%} \\
\hline
\end{tabular}
\caption{Class Distribution in the Medicare Provider-Level Dataset}
\label{tab:class-dist}
\end{table}

The class distribution confirms a significant class imbalance: approximately
9.35\% of providers carry the ``Potential Fraud'' label. This imbalance reflects
the real-world distribution of fraud in healthcare claims data and necessitates
specific measures such as class-weight balancing and appropriate evaluation
metrics beyond accuracy alone.

\section{Data Preprocessing Results}

\subsection{Provider-Level Feature Engineering}

The raw Medicare claim records were aggregated at the provider level using
ten engineered features. The feature engineering step processed all 558,211
combined inpatient and outpatient claims to produce the ten provider-level
summary statistics described in Chapter~\ref{ch:design}. Missing values arising
from providers with no inpatient or no outpatient claims were filled with zero.

\begin{table}[h!]
\centering
\renewcommand{\arraystretch}{1.3}
\begin{tabular}{|l|r|}
\hline
\textbf{Feature} & \textbf{Description} \\
\hline
TotalClaims & Total claims submitted (all types) \\
\hline
TotalReimbursement & Total InscClaimAmtReimbursed (USD) \\
\hline
AverageReimbursement & Mean reimbursement per claim (USD) \\
\hline
UniquePatients & Distinct beneficiary count (BeneID) \\
\hline
InpatientClaims & Count of inpatient claim records \\
\hline
InpatientReimbursement & Total inpatient reimbursement (USD) \\
\hline
InpatientAvgReimbursement & Mean inpatient reimbursement per claim (USD) \\
\hline
OutpatientClaims & Count of outpatient claim records \\
\hline
OutpatientReimbursement & Total outpatient reimbursement (USD) \\
\hline
OutpatientAvgReimbursement & Mean outpatient reimbursement per claim (USD) \\
\hline
\end{tabular}
\caption{Ten Provider-Level Engineered Features}
\label{tab:features-result}
\end{table}

The resulting provider-level feature matrix had dimensions of 5,410 rows and
10 columns. All feature values were of floating-point type. The provider with the
highest total reimbursement and the highest inpatient reimbursement were also among
the providers assigned a High risk tier by the Random Forest model.

\subsection{Claim-Level Screening Flag Distribution}

Five binary screening flags were computed for all 558,211 claims. The distribution
of the total number of screening flags per claim (ScreeningFlagsCount) provides an
overview of the claim-level anomaly landscape in the dataset.

\begin{table}[h!]
\centering
\renewcommand{\arraystretch}{1.3}
\begin{tabular}{|l|l|}
\hline
\textbf{Screening Rule} & \textbf{Trigger Condition} \\
\hline
Reimbursement Outlier & InscClaimAmtReimbursed $>$ 99th percentile for claim type \\
\hline
Duplicate Billing & Same BeneID + Provider + ClaimStartDt on multiple claims \\
\hline
Duration Anomaly & ClaimDurationDays $<$ 0 or LengthOfStayDays $<$ 0 \\
\hline
Complex Procedure & Outpatient claim with $\geq 3$ procedure codes \\
\hline
High-Frequency Beneficiary & Beneficiary with $\geq 15$ claims in the dataset window \\
\hline
\end{tabular}
\caption{Five Claim-Level Transparent Screening Rules}
\label{tab:screening-rules}
\end{table}

The AI Review Recommendation field is derived from the combination of the provider
risk level and the screening flag count. Claims associated with High-risk providers
and a non-zero reimbursement are assigned the ``Manual Investigation'' recommendation.
Claims associated with High-risk providers or with two or more flags receive a
``Risk Review'' recommendation. Claims with one flag or from Medium-risk providers
receive a ``Priority Review'' recommendation. All remaining claims receive a
``Routine Review'' recommendation.

\subsection{Train-Test Split}

The provider-level dataset was divided into training and testing subsets using an
80:20 stratified split, implemented using \verb|StratifiedKFold| from scikit-learn.
Stratification ensures that the proportion of potentially fraudulent providers
(9.35\%) is preserved in both the training and testing partitions.

\begin{table}[h!]
\centering
\renewcommand{\arraystretch}{1.3}
\begin{tabular}{|l|r|r|r|}
\hline
\textbf{Partition} & \textbf{Total Providers} & \textbf{Fraud (Yes)} & \textbf{Non-Fraud (No)} \\
\hline
Training Set & 4,328 & $\approx$405 & $\approx$3,923 \\
\hline
Test Set & 1,082 & $\approx$101 & $\approx$981 \\
\hline
\textbf{Total} & \textbf{5,410} & \textbf{506} & \textbf{4,904} \\
\hline
\end{tabular}
\caption{Stratified 80:20 Train-Test Split for Provider-Level Dataset}
\label{tab:split}
\end{table}

\section{NFHS-5 Medical Validation Results}

\subsection{Nationwide Benchmark Results}

The six clinical consistency rules were applied to the complete NFHS-5 dataset of
724,115 records. The rules identify medically inconsistent combinations of
self-reported diagnosis, treatment, medication, and glucose measurement. A weighted
risk score was computed for each record as:

\[
\text{RiskScore} = 3 R_{1A} + 3 R_{1B} + 5 R_{2A} + 5 R_{2B} + 4 R_3 + 2 R_4
\]

where $R_{1A}$ is the diabetes-without-treatment indicator, $R_{1B}$ is the
hypertension-without-treatment indicator, $R_{2A}$ is the treatment-without-diabetes
indicator, $R_{2B}$ is the treatment-without-hypertension indicator, $R_3$ is the
high-glucose-without-diagnosis indicator, and $R_4$ is the
diabetes-without-medication indicator. Each $R_i \in \{0, 1\}$.

Records were classified into five risk categories based on their cumulative risk
score. The risk category thresholds and the resulting distribution across all
724,115 records are shown in Table~\ref{tab:nfhs-nationwide}.

\begin{table}[h!]
\centering
\renewcommand{\arraystretch}{1.3}
\begin{tabular}{|l|c|r|r|}
\hline
\textbf{Risk Category} & \textbf{Score Threshold} & \textbf{Records} & \textbf{Percentage} \\
\hline
Clean & $= 0$ & 703,601 & 97.17\% \\
\hline
Low & $\leq 2$ & 5,203 & 0.72\% \\
\hline
Medium & $\leq 4$ & 12,512 & 1.73\% \\
\hline
High & $\leq 6$ & 2,328 & 0.32\% \\
\hline
Critical & $> 6$ & 471 & 0.07\% \\
\hline
\textbf{Total Flagged} & -- & \textbf{20,514} & \textbf{2.83\%} \\
\hline
\textbf{Total} & -- & \textbf{724,115} & \textbf{100.00\%} \\
\hline
\end{tabular}
\caption{NFHS-5 Nationwide Medical Validation Results (724,115 Records)}
\label{tab:nfhs-nationwide}
\end{table}

The results show that the overwhelming majority of NFHS-5 records (97.17\%) did not
trigger any of the six clinical inconsistency rules and were classified as Clean.
A total of 20,514 records (2.83\%) were flagged with at least one inconsistency.
Among flagged records, 2,799 were categorised as High or Critical risk,
representing individuals with the most clinically inconsistent combinations of
reported health information.

\subsection{Rule-Wise Validation Results}

The frequency of each individual rule trigger across the 724,115 records provides
additional insight into the nature of the identified inconsistencies.

\begin{table}[h!]
\centering
\renewcommand{\arraystretch}{1.3}
\begin{tabular}{|l|c|r|r|}
\hline
\textbf{Rule} & \textbf{Weight} & \textbf{Triggered Count} & \textbf{Percentage} \\
\hline
R1A: Diabetes without treatment & 3 & 2,679 & 0.37\% \\
\hline
R1B: Hypertension without treatment & 3 & 8,340 & 1.15\% \\
\hline
R2A: Treatment without diabetes diagnosis & 5 & 0 & 0.00\% \\
\hline
R2B: Treatment without hypertension diagnosis & 5 & 0 & 0.00\% \\
\hline
R3: High glucose ($\geq$200 mg/dL) without diagnosis & 4 & 4,790 & 0.66\% \\
\hline
R4: Diabetes without prescribed medication & 2 & 7,894 & 1.09\% \\
\hline
\end{tabular}
\caption{Rule-Wise NFHS-5 Medical Validation Trigger Frequencies (724,115 Records)}
\label{tab:nfhs-rules}
\end{table}

The most frequently triggered rule is R1B (Hypertension without treatment), affecting
8,340 records (1.15\%), followed by R4 (Diabetes without prescribed medication)
with 7,894 records (1.09\%) and R3 (High blood glucose without a diabetes diagnosis)
with 4,790 records (0.66\%). Rules R2A and R2B (treatment reported without a
corresponding diagnosis) triggered zero times in this dataset. This finding is
consistent with the survey design, wherein respondents who reported receiving
treatment were also expected to report the corresponding diagnosis. The absence of
R2A and R2B triggers therefore reflects internal consistency in the survey response
patterns rather than an absence of the rule logic.

\subsection{NFHS-5 Screened Sample Results}

A 35,000-record screened sample was extracted from the NFHS-5 data for integration
into the interactive dashboard. Of these 35,000 records, 1,213 (3.46\%) were
identified with at least one rule trigger. The risk category distribution within the
sample is shown below.

\begin{table}[h!]
\centering
\renewcommand{\arraystretch}{1.3}
\begin{tabular}{|l|r|r|}
\hline
\textbf{Risk Category} & \textbf{Sample Records} & \textbf{Percentage} \\
\hline
Clean & 33,787 & 96.53\% \\
\hline
Low & 484 & 1.38\% \\
\hline
Medium & 567 & 1.62\% \\
\hline
High & 150 & 0.43\% \\
\hline
Critical & 12 & 0.03\% \\
\hline
\textbf{Total} & \textbf{35,000} & \textbf{100.00\%} \\
\hline
\end{tabular}
\caption{NFHS-5 Screened Sample Risk Category Distribution (35,000 Records)}
\label{tab:nfhs-sample}
\end{table}

The sample rule-wise trigger counts (R1A: 147, R1B: 460, R2A: 0, R2B: 0, R3: 137,
R4: 640) are consistent in relative proportions with the nationwide benchmark
statistics, confirming that the sample is representative of the broader NFHS-5
population's clinical inconsistency profile.

\subsection{Interpretation of NFHS-5 Results}

The NFHS-5 validation component does not produce confirmed fraud findings.
The identified inconsistencies represent medically unusual combinations of
self-reported survey responses and clinical measurements. These patterns can arise
from reporting errors, incomplete survey participation, data entry discrepancies,
or genuine clinical inconsistencies that may warrant further follow-up. The
proposed system treats these records as candidates for review rather than
as confirmed cases.

The concentration of inconsistencies in the Hypertension-without-treatment (R1B)
and Diabetes-without-medication (R4) categories suggests that awareness or
treatment gaps for chronic conditions may be the most prevalent source of clinically
inconsistent patterns in the NFHS-5 data. The high-glucose-without-diagnosis (R3)
category further highlights the presence of respondents with elevated glucose
measurements who had not reported a formal diabetes diagnosis, which could indicate
undiagnosed diabetes among survey respondents.

\section{Medicare Provider-Level Fraud Detection Results}

\subsection{Cross-Validation Model Evaluation}

Three supervised machine learning models were trained and evaluated using five-fold
Stratified Cross-Validation. The cross-validation was performed on the full
provider-level dataset of 5,410 providers, maintaining the class distribution in
each fold. Six performance metrics were computed for each model.

The cross-validation results from the authoritative \verb|final_model_comparison.csv|
are presented in Table~\ref{tab:cv-results}.

\begin{table}[h!]
\centering
\renewcommand{\arraystretch}{1.4}
\begin{tabular}{|l|c|c|c|c|c|c|}
\hline
\textbf{Model} & \textbf{Accuracy} & \textbf{Precision} & \textbf{Recall} &
\textbf{F1-Score} & \textbf{ROC-AUC} & \textbf{Avg Precision} \\
\hline
Random Forest & 92.98\% & 68.93\% & 44.87\% & 54.19\% & 92.33\% & 66.11\% \\
\hline
Logistic Regression & 87.39\% & 41.40\% & 83.40\% & 55.32\% & 93.53\% & 68.67\% \\
\hline
XGBoost & 88.82\% & 44.73\% & 80.83\% & 57.56\% & 93.55\% & 69.00\% \\
\hline
\end{tabular}
\caption{Five-Fold Stratified Cross-Validation Results for Provider-Level Fraud Detection}
\label{tab:cv-results}
\end{table}

The results reveal distinctly different precision-recall characteristics across the
three models. Random Forest achieves the highest accuracy (92.98\%) and the highest
precision (68.93\%), meaning that when it classifies a provider as potentially
fraudulent, it is correct approximately 68.93\% of the time. However, its recall
is comparatively low (44.87\%), indicating that it misses a substantial portion of
the truly fraudulent providers in the dataset.

Logistic Regression and XGBoost both achieve significantly higher recall values
(83.40\% and 80.83\% respectively), at the cost of lower precision (41.40\% and
44.73\%). These models identify a larger fraction of the potentially fraudulent
providers but also generate more false positive flags. This precision-recall
trade-off is an inherent characteristic of imbalanced classification problems and
is deliberately surfaced to the investigator rather than resolved by choosing a
single ``best'' model.

The ROC-AUC scores are high across all three models (92.33\% for Random Forest,
93.53\% for Logistic Regression, 93.55\% for XGBoost), indicating strong
discriminative ability at the ranking level. Average Precision scores (66.11\%,
68.67\%, 69.00\%) provide a more conservative estimate of precision at all recall
levels and are more appropriate for imbalanced classification than the area under
the standard ROC curve.

\subsection{Model Metadata Cross-Reference}

The project contains two recorded sets of five-fold CV metrics: one in
\verb|final_model_comparison.csv| and one in \verb|model_metadata.json|.
The two sets differ slightly for Random Forest and XGBoost, reflecting different
training configurations across development runs. The Random Forest metadata
reports Accuracy = 93.03\%, Precision = 68.00\%, Recall = 47.64\%, F1 = 55.94\%,
ROC-AUC = 92.26\%, Average Precision = 66.68\%. The XGBoost metadata reports
Accuracy = 88.65\%, Precision = 44.26\%, Recall = 80.83\%, F1 = 57.17\%,
ROC-AUC = 93.33\%, Average Precision = 68.45\%. The Logistic Regression metrics
are identical across both sources. The dashboard's ML Performance page displays
both sources and their per-metric differences to ensure transparency.

\subsection{Threshold Analysis}

A key feature of the system is the ability to adjust the classification threshold
interactively via the ML Model Performance dashboard page. The threshold slider
operates on the held-out evaluation predictions stored in
\verb|eval_predictions.parquet|, which were generated by
\verb|scripts/generate_eval_data.py| using an 80:20 stratified split (4,328 train,
1,082 test providers).

The default operational thresholds are:
\begin{itemize}
\item \textbf{Default Fraud Threshold} (\verb|DEFAULT_FRAUD_THRESHOLD|): 0.30 — providers with a predicted fraud probability $\geq$ 0.30 are assigned the Medium or High risk tier.
\item \textbf{High Risk Threshold} (\verb|DEFAULT_HIGH_RISK_THRESHOLD|): 0.70 — providers with a predicted fraud probability $\geq$ 0.70 are assigned the High risk tier and trigger ``Manual Investigation'' for all associated claims.
\end{itemize}

Lowering the threshold increases recall (more potentially fraudulent providers
are flagged) while reducing precision (more non-fraudulent providers are
incorrectly flagged). Raising the threshold has the opposite effect. The system
supports threshold values in the range $[0.0, 1.0]$ and dynamically recomputes
accuracy, precision, recall, F1-score, and the confusion matrix counts from the
holdout predictions without retraining. ROC-AUC and Average Precision are
threshold-independent metrics derived from the complete score distribution.

\subsection{Provider Risk Tier Distribution}

After applying the default thresholds to the Random Forest model predictions over
the full 5,410-provider dataset, the risk tier distribution is as follows:

\begin{table}[h!]
\centering
\renewcommand{\arraystretch}{1.3}
\begin{tabular}{|l|c|r|r|}
\hline
\textbf{Risk Tier} & \textbf{Probability Range} & \textbf{Providers} & \textbf{Percentage} \\
\hline
Low & $< 0.30$ & 4,882 & 90.24\% \\
\hline
Medium & $0.30 \leq p < 0.70$ & 133 & 2.46\% \\
\hline
High & $\geq 0.70$ & 395 & 7.30\% \\
\hline
\textbf{Total} & -- & \textbf{5,410} & \textbf{100.00\%} \\
\hline
\end{tabular}
\caption{Provider Risk Tier Distribution under Default Thresholds (RF Model)}
\label{tab:risk-tiers}
\end{table}

The High-risk tier contains 395 providers, compared to the 506 providers with the
``Yes'' PotentialFraud label. The discrepancy reflects the conservative nature of
the 0.70 high-risk threshold. Many providers labelled as potentially fraudulent in
the ground-truth data are assigned to the Medium tier under the default threshold,
consistent with the cross-validation recall of 44.87\%.

\section{Feature Importance Analysis}

\subsection{Random Forest Gini Feature Importances}

The Random Forest model's feature importances, computed from the mean decrease in
Gini impurity across all 100 decision trees, are reported in Table~\ref{tab:feat-imp}.
These importances reflect the contribution of each feature to the model's
classification performance across the training data.

\begin{table}[h!]
\centering
\renewcommand{\arraystretch}{1.3}
\begin{tabular}{|l|r|}
\hline
\textbf{Feature} & \textbf{Gini Importance} \\
\hline
TotalReimbursement & 22.98\% \\
\hline
InpatientReimbursement & 16.44\% \\
\hline
InpatientClaims & 11.31\% \\
\hline
InpatientAvgReimbursement & 8.91\% \\
\hline
UniquePatients & 7.58\% \\
\hline
AverageReimbursement & 7.33\% \\
\hline
TotalClaims & 7.22\% \\
\hline
OutpatientReimbursement & 6.83\% \\
\hline
OutpatientClaims & 5.91\% \\
\hline
OutpatientAvgReimbursement & 5.49\% \\
\hline
\end{tabular}
\caption{Random Forest Gini Feature Importances (from model\_metadata.json)}
\label{tab:feat-imp}
\end{table}

The results indicate that reimbursement-related features are the most influential
predictors in the Random Forest model. \verb|TotalReimbursement| is the single most
important feature with a Gini importance of 22.98\%, followed by
\verb|InpatientReimbursement| (16.44\%) and \verb|InpatientClaims| (11.31\%).
Together, the top three features account for over 50\% of the model's discriminative
power. Outpatient features (OutpatientClaims at 5.91\% and OutpatientAvgReimbursement
at 5.49\%) have the lowest individual importance values, suggesting that inpatient
billing behaviour is the stronger differentiator between potentially fraudulent and
non-fraudulent providers in this dataset.

This finding is consistent with known healthcare fraud patterns, where fraudulent
providers may inflate inpatient billing volumes and reimbursement claims, submit
charges for services not rendered, or falsely escalate the complexity of reported
procedures to maximise reimbursement from the Medicare programme.

\subsection{SHAP-Based Local Explanations}

In addition to global Gini importances, the system provides local SHAP (SHapley
Additive exPlanations) feature attribution for individual provider predictions. SHAP
TreeExplainer is applied to the trained Random Forest model using the
\verb|explain_provider_prediction()| function in \verb|src/explainability.py|.

For each provider, SHAP computes a contribution value $\phi_i$ for every feature $i$
such that:

\[
f(x) = \phi_0 + \sum_{i=1}^{10} \phi_i
\]

where $f(x)$ is the model output (fraud probability) for provider $x$, $\phi_0$ is
the base value (expected model output over the training distribution), and $\phi_i$
is the contribution of feature $i$ to the difference between the provider's
predicted probability and the base value.

Features with positive SHAP values push the predicted probability higher (towards
the fraud class), while features with negative SHAP values reduce the predicted
probability. The system displays the top three risk-elevating features and the top
two risk-reducing features as plain-English investigator findings, enhancing
transparency for non-technical users.

For example, a provider with elevated \verb|TotalReimbursement| and
\verb|InpatientReimbursement| values substantially above the provider population
mean will receive large positive SHAP contributions for those features, pushing the
predicted fraud probability upward. The Provider Risk Analysis dashboard page renders
these contributions as a horizontal waterfall chart using the
\verb|plot_shap_contributions()| function in \verb|src/visualizations.py|.

\section{Interactive Dashboard Validation}

\subsection{Claim Search Performance}

The DuckDB analytical database (\verb|claims_fraud.duckdb|, 85\,MB) supports
sub-second queries across 558,211 claim records. Indexed lookups by \verb|ClaimID|,
\verb|BeneID|, and \verb|Provider| consistently complete in under 10 milliseconds
on the test hardware. The \verb|search_claims()| function in
\verb|src/claims_repository.py| uses parameterised DuckDB queries with an explicit
column whitelist to prevent SQL injection, and supports pagination with a configurable
limit (maximum 500 records per page).

\subsection{Investigation Workspace Functionality}

The SQLite-backed investigation workspace (\verb|investigation_workspace.db|)
successfully maintained transactional consistency across test cases covering case
creation, note addition, status transitions, and audit log recording. The
\verb|create_or_get_case()| function uses an \verb|ON CONFLICT DO NOTHING| strategy
to guarantee exactly one case per investigation target (identified by the combination
of \verb|target_type| and \verb|target_id|), preventing duplicate case creation
even under concurrent access patterns.

\subsection{Automated Test Suite Results}

The automated test suite consists of five test modules and is executed using pytest.
The suite covers data integrity, DuckDB claim queries, model inference, screening
logic, and investigation workspace persistence. All tests passed on the verified
production dataset. The test modules and their scope are summarised as follows:

\begin{table}[h!]
\centering
\renewcommand{\arraystretch}{1.3}
\begin{tabular}{|l|l|}
\hline
\textbf{Test Module} & \textbf{Scope} \\
\hline
\verb|test_data_validation.py| & Parquet file integrity, DuckDB existence, 10 feature columns \\
\hline
\verb|test_claim_queries.py| & DuckDB lookups (CLM46614), non-existent IDs, provider profiles \\
\hline
\verb|test_predictions.py| & Model loading, $0.0 \leq p \leq 1.0$, risk tiers, SHAP values \\
\hline
\verb|test_dashboard_logic.py| & SQLite transactions, audit logs, screening rules, NFHS loading \\
\hline
\verb|test_investigation_dashboard.py| & Date filters and stored screening evidence \\
\hline
\end{tabular}
\caption{Automated Test Suite Modules and Their Coverage}
\label{tab:tests}
\end{table}

\section{Discussion of Results}

The experimental results demonstrate that the proposed two-component framework
is technically feasible and produces meaningful outputs for the stated objectives.

The NFHS-5 medical validation component successfully identifies a small but
non-negligible subset of records (2.83\% of 724,115 records) with clinical
inconsistencies, including the well-known patterns of untreated hypertension and
unmedicated diabetes, as well as elevated blood glucose without a recorded diagnosis.
The rule-based approach is transparent and interpretable, and its outputs are well
suited for clinical follow-up and healthcare programme quality assurance.

The Medicare provider-level component demonstrates that aggregate billing features
can meaningfully distinguish providers labelled as potentially fraudulent from those
labelled as non-fraudulent. The high ROC-AUC values (approximately 92--94\%) across
all three models indicate strong discriminative ability, even in the presence of
significant class imbalance. The divergence in precision-recall characteristics
across models is a useful property of the system, because different investigation
contexts may prioritise either catching more suspicious providers (favouring XGBoost
or Logistic Regression for their higher recall) or minimising unnecessary
investigations (favouring Random Forest for its higher precision).

The SHAP-based explainability layer adds interpretability to the provider-level
predictions. By identifying the specific features that drive a provider's elevated
risk score, the system provides actionable investigation leads rather than an
opaque binary output. This is particularly important in a healthcare fraud
investigation context, where model predictions must be scrutinised alongside source
evidence before any enforcement or billing action is taken.

\section{Limitations}

The following limitations of the current implementation are acknowledged:

\begin{enumerate}

\item \textbf{Provider-Level Ground Truth Only:} The Medicare dataset contains ground
truth labels only at the provider level (\verb|PotentialFraud| Yes/No). Individual
claim-level fraud labels are not available. Consequently, claim-level machine
learning models cannot be trained. The claim-level screening rules are heuristic and
transparent indicators, not supervised model outputs.

\item \textbf{Synthetic Medicare Data:} The CMS SynPUF dataset is a synthetic
representation of Medicare data. Results obtained from this dataset may not fully
generalise to real Medicare claims with the same accuracy, as synthetic data
generation may not preserve all fraud-indicative distributional properties.

\item \textbf{NFHS-5 Full Benchmark Verification:} The nationwide benchmark summary
figures (724,115 records, 20,514 flagged) are recorded in
\verb|data/nfhs/nfhs5_rule_summary.json|, but the full raw NFHS-5 benchmark dataset
is not present in the current project folder. These figures therefore cannot be
independently recomputed from within the current project artefacts, though they are
consistent with the proportional results observed in the 35,000-record sample.

\item \textbf{Class Imbalance:} Despite class-weight balancing in Logistic Regression
and the use of \verb|scale_pos_weight| in XGBoost, the class imbalance (9.35\%
positive class) remains a challenge. The Random Forest model's recall of 44.87\%
means that approximately 55\% of truly fraudulent providers are not flagged at the
default threshold.

\item \textbf{No Real-Time Streaming:} The current system operates in batch mode on
a static dataset. It does not support real-time ingestion of incoming Medicare claims
or continuous retraining of models as new labelled data becomes available.

\item \textbf{Provider-Scope of SHAP:} SHAP attributions in the current system are
computed at the provider level using aggregated features. Individual claim-level
explanations using SHAP are not currently supported.

\end{enumerate}

